% =====================================================================
%  preamble.tex - 패키지, 폰트, 스타일, semantic 매크로
%  집필 규칙: ../AUTHORING.md
%
%  엔진은 XeLaTeX 고정. 로컬 TeX Live 는 basic 스킴이라 kotex, tcolorbox,
%  mdframed, framed, siunitx, biblatex 가 없다. 한글은 xeCJK 로 처리한다.
% =====================================================================

% ---------------------------------------------------------------- 한글
\usepackage{fontspec}
\usepackage{xeCJK}

% CJKspace=true 는 필수다. 없으면 한글 단어 사이 띄어쓰기가 사라진다.
\xeCJKsetup{CJKspace=true}

% 한글 폰트에는 이탤릭 자형이 없다. \analogy{} 가 그냥 무시되지 않도록
% 기울임을 합성(FakeSlant)한다.
\setCJKmainfont{Apple SD Gothic Neo}[
  BoldFont       = {Apple SD Gothic Neo Bold},
  ItalicFont     = {Apple SD Gothic Neo},
  ItalicFeatures = {FakeSlant = 0.167},
]
\setCJKsansfont{Apple SD Gothic Neo}

% Latin Modern 은 fontconfig 에 이름이 등록되어 있지 않다(basic 스킴).
% 이름이 아니라 파일명으로 지정해야 XeLaTeX 이 찾는다.
\setmainfont{lmroman10-regular.otf}[
  BoldFont       = lmroman10-bold.otf,
  ItalicFont     = lmroman10-italic.otf,
  BoldItalicFont = lmroman10-bolditalic.otf,
]
\setmonofont{Menlo}[Scale = 0.86]

% 한글은 어절 단위로 끊어지므로 하이픈을 끄고 자간 신축을 허용한다.
\XeTeXlinebreaklocale ``ko''
\XeTeXlinebreakskip = 0pt plus 0.1em minus 0.01em

% ---------------------------------------------------------------- 수식
\usepackage{amsmath}
\usepackage{amssymb}
\usepackage{amsthm}

% ------------------------------------------------------- 표, 그림, 코드
\usepackage{graphicx}
\usepackage{booktabs}
\usepackage{tabularx}
\usepackage{array}
\usepackage[font=small, labelfont=bf]{caption}
\usepackage{subcaption}
\usepackage{listings}
\usepackage{tikz}
\usetikzlibrary{arrows.meta, positioning, fit, backgrounds, calc}

% -------------------------------------------------------------- 레이아웃
\usepackage[a4paper, margin=2.6cm, headheight=14pt]{geometry}
\usepackage{fancyhdr}
\usepackage{titlesec}
\usepackage{enumitem}
\usepackage{xcolor}
\usepackage{microtype}

% ------------------------------------------------- 상호참조 (마지막 로드)
\usepackage{natbib}
\usepackage[colorlinks = true, linkcolor = bookaccent,
            citecolor = bookaccent, urlcolor = bookaccent]{hyperref}
\usepackage{bookmark}
\usepackage[nameinlink]{cleveref}   % hyperref 뒤에 와야 한다

% =====================================================================
%  색
%  hw1 그림들과 같은 팔레트를 쓴다. 책과 그림의 색이 어긋나면 독자가
%  둘을 다른 출처로 인식한다.
% =====================================================================
\definecolor{bookaccent}{HTML}{2A78D6}
\definecolor{bookwarn}{HTML}{EB6834}
\definecolor{bookink}{HTML}{0B0B0B}
\definecolor{bookmuted}{HTML}{52514E}
\definecolor{bookrule}{HTML}{D9D8D4}
\definecolor{booksurface}{HTML}{F7F7F5}
\definecolor{bookwarnbg}{HTML}{FDF1EC}
\definecolor{bookaccentbg}{HTML}{EDF3FC}

% =====================================================================
%  Semantic 매크로 (AUTHORING.md 5절)
%  의미 단위로 표시한다. 스타일은 여기 한 곳에서만 바꾼다.
% =====================================================================
\newcommand{\term}[1]{\textbf{#1}}          % 정의되는 개념
\newcommand{\analogy}[1]{\textit{#1}}       % 비유 표현
\newcommand{\contains}{\ensuremath{\supset}}% 담김/포함 관계

% 자주 쓰는 수식 기호
\newcommand{\E}{\operatorname{E}}
\newcommand{\Var}{\operatorname{Var}}
\newcommand{\IQR}{\mathrm{IQR}}
\newcommand{\sd}{\mathrm{SD}}

% =====================================================================
%  콜아웃 상자
%  tcolorbox 가 없으므로 \fcolorbox 로 만든다. lrbox 에 담아 두었다가
%  통째로 그리므로 페이지를 넘지 않는다 - 짧게 쓸 것.
% =====================================================================
\newsavebox{\bookcallout}
\newlength{\bookcalloutwd}

% \newenvironment 의 종료부에서는 #1, #2 를 쓸 수 없다(인자는 시작부에만
% 전달된다). 그래서 색 이름을 매크로에 담아 두었다가 종료부에서 꺼내 쓴다.
\newcommand{\calloutborder}{bookaccent}
\newcommand{\calloutbg}{bookaccentbg}

\newenvironment{calloutbox}[3]%
  % #1 테두리색  #2 배경색  #3 머리말
  {\renewcommand{\calloutborder}{#1}\renewcommand{\calloutbg}{#2}%
   \setlength{\bookcalloutwd}{\dimexpr\linewidth-2\fboxsep-2\fboxrule\relax}%
   \par\medskip\noindent
   \begin{lrbox}{\bookcallout}\begin{minipage}{\bookcalloutwd}%
   \setlength{\parskip}{4pt}%
   \textbf{\color{#1}#3}\par\nobreak}%
  {\end{minipage}\end{lrbox}%
   \setlength{\fboxrule}{1pt}%
   \fcolorbox{\calloutborder}{\calloutbg}{\usebox{\bookcallout}}\par\medskip}

% 핵심: 그 절에서 반드시 남겨야 할 한 가지
\newenvironment{keyidea}[1][핵심]%
  {\begin{calloutbox}{bookaccent}{bookaccentbg}{#1}\color{bookink}}%
  {\end{calloutbox}}

% 함정: 내가 실제로 틀렸던 것, 또는 틀리기 쉬운 것
\newenvironment{pitfall}[1][함정]%
  {\begin{calloutbox}{bookwarn}{bookwarnbg}{#1}\color{bookink}}%
  {\end{calloutbox}}

% 여담: 본문 흐름에서 빠져도 되는 배경 지식
\newenvironment{aside}[1][참고]%
  {\par\medskip\small\color{bookmuted}\noindent\textbf{#1}\quad\ignorespaces}%
  {\par\normalsize\color{bookink}\medskip\ignorespacesafterend}

% =====================================================================
%  정리 환경
% =====================================================================
\theoremstyle{definition}
\newtheorem{definition}{정의}[chapter]
\newtheorem{example}{예제}[chapter]
\theoremstyle{plain}
\newtheorem{theorem}{정리}[chapter]

\crefname{definition}{정의}{정의}
\crefname{example}{예제}{예제}
\crefname{chapter}{장}{장}
\crefname{section}{절}{절}
\crefname{figure}{그림}{그림}
\crefname{table}{표}{표}

% =====================================================================
%  시각 스타일
% =====================================================================
\pagestyle{fancy}
\fancyhf{}
\fancyhead[L]{\small\color{bookmuted}\nouppercase{\leftmark}}
\fancyhead[R]{\small\color{bookmuted}\thepage}
\renewcommand{\headrulewidth}{0.4pt}
\renewcommand{\headrule}{\color{bookrule}\hrule height 0.4pt}
\fancypagestyle{plain}{\fancyhf{}%
  \fancyhead[R]{\small\color{bookmuted}\thepage}%
  \renewcommand{\headrulewidth}{0pt}}

\titleformat{\chapter}[display]
  {\normalfont\huge\bfseries\color{bookink}}
  {\color{bookaccent}\large\chaptertitlename\ \thechapter}{8pt}{\Huge}
\titlespacing*{\chapter}{0pt}{-20pt}{28pt}
\titleformat{\section}{\normalfont\Large\bfseries\color{bookink}}
  {\color{bookaccent}\thesection}{0.7em}{}
\titleformat{\subsection}{\normalfont\large\bfseries\color{bookink}}
  {\color{bookaccent}\thesubsection}{0.7em}{}

% ``In brief.'' 블록의 description 라벨을 고정폭으로 세운다 (AUTHORING.md 1.6).
\newenvironment{inbrief}%
  {\paragraph{In brief.}%
   \begin{description}[leftmargin=5.6em, labelsep=0.8em, labelwidth=4.4em,
     align=left, font={\bfseries\color{bookaccent}}, itemsep=3pt, topsep=4pt,
     parsep=0pt]}%
  {\end{description}\medskip}

\lstdefinestyle{book}{
  basicstyle       = \ttfamily\small,
  keywordstyle     = \color{bookaccent},
  commentstyle     = \color{bookmuted}\itshape,
  stringstyle      = \color{bookwarn},
  backgroundcolor  = \color{booksurface},
  frame            = single,
  rulecolor        = \color{bookrule},
  breaklines       = true,
  showstringspaces = false,
  columns          = fullflexible,
  xleftmargin      = 10pt,
  framexleftmargin = 10pt,
  extendedchars    = true,
  inputencoding    = utf8,
}
\lstset{style=book}

\setlist{itemsep=2pt, topsep=4pt, parsep=0pt}
\setlength{\parindent}{0pt}
\setlength{\parskip}{6pt}
